% AASTeX v7.0.1 has no RNAAS class option (it existed in v6.3.x and was dropped),
% so this is the standard manuscript style. Confirm the current Research Notes
% template with AAS before submitting; the journal typesets RNAAS itself.
\documentclass{aastex701}

\newcommand{\flux}{\mathrm{molecules\ cm^{-2}\ s^{-1}}}
\newcommand{\Fch}{F_{\mathrm{CH_4}}}

\shorttitle{The abiotic methane ceiling}
\shortauthors{Sharma}

\begin{document}

\title{The Abiotic Methane Ceiling Sets a Floor on Detectable Biospheres}

% Add your ORCID as \author[0000-...]{...} if you have one.
\author{Siddhant Sharma}
\affiliation{Blue Marble Space Institute of Science, Seattle, WA 98154, USA}
\email{siddhant.sharma@bmsis.org}

\begin{abstract}
Flux based biosignature assessment compares an inferred surface methane flux against
a ceiling on abiotic production. The ceiling in current use is a deliberately
conservative upper bound, adopted as a single number. It stands above every other published maximum for the same
process, by factors of four to fifty-five, above anything its own generating equation
yields at modern Earth crustal production rates, and a factor of three below Earth's
own biological methane flux. The choice
of ceiling therefore fixes the smallest biosphere an observation can reveal, and the
span of defensible choices is wider than the spectroscopic uncertainty of the
retrieval it is compared against.
\end{abstract}

\section{A bound doing the work of a likelihood} \label{sec:setup}

\citet{wogan2026} infer surface gas fluxes for an Archean-like TRAPPIST-1\,e by
inverting a coupled photochemical and climate model, then ask whether the inferred
fluxes are consistent with life. The step from flux to biosignature runs through a
set definition: methane fluxes above $3.7\times10^{10}\ \flux$ are taken to exceed
what ocean bottom serpentinization can supply. That value is the 99.7th percentile
of the distribution of maximum abiotic fluxes computed by \citet{kt2018}, adopted,
as \citet{wogan2026} say, for simplicity.

As a set boundary the choice is sound and conservative by design. Its consequences
appear at the following step, where the complement of the biosignature region is
taken as the abiotic alternative and the resulting odds ratio is used as a proxy for
the Bayes factor between life and no life. A bound chosen to be generous then does
the work of a likelihood.

The construction carries a second consequence. The oxygen biosignature is defined as
the subset of $\mathcal{H}_{\mathrm{CH_4}}$ with a net positive oxygen flux, and the
life-consistent region $\mathcal{L}$ as its union with the methane biosignature,
which is $\mathcal{H}_{\mathrm{CH_4}}$ with the reducing volcanic false positive
$\mathcal{V}$ removed. The union therefore restores part of $\mathcal{V}$: their
Table 2 masses give $P(\mathcal{V}\cap\mathcal{B}_{\mathrm{O_2}}) = P(\mathcal{L})
- P(\mathcal{H}_{\mathrm{CH_4}}) + P(\mathcal{V})$, at least half of $\mathcal{V}$
in both prior and posterior once two-decimal rounding is propagated. What one branch
excludes as a false positive, the other partly readmits.

\section{Locating the ceiling} \label{sec:locating}

Three comparisons place it, and the first needs no view on what Earth actually
produces, since it sets the ceiling against other published maxima.
\citet{guzman2013} modelled methane
production by serpentinization from the kinetics of the reactions involved, with
carbon dioxide as the limiting reactant, and report a maximum surface flux of
$6.8\times10^{8}\ \flux$ for an Earth mass planet, 1.74 dex below the ceiling.
\citet{merdith2020} give an upper estimate in which the methane flux approaches the
hydrogen flux, 1.12 dex below it. Even the Archean maximum of \citet{fiebig2007},
which describes a more vigorous Earth than the one being modelled, falls 0.60 dex
short. No published maximum for this process reaches the adopted ceiling, and the
three fall short of it by factors of four, thirteen and fifty-five.

The second uses the generating equation itself. The lower corner of the four
parameter ranges sampled by \citet{kt2018} is built entirely from Earth calibrated
values: their modern crustal production rate, 10 wt\% FeO in newly produced crust,
the 0.25\% conversion of FeO to hydrogen that they calibrate against observation,
and the hydrogen to methane ratio of 12 measured in ultramafic hosted vent fluids.
Their Equation 6 then returns 0.012 Tmol yr$^{-1}$, or $\log \Fch = 7.67$, a little
below the directly measured mid-ocean ridge values. Their distribution runs from
what Earth does to the most generous case they were willing to entertain, and the
adopted percentile sits 2.90 dex above its lower end.

The third asks what that equation reaches without leaving Earth-like conditions.
Holding crustal production at the modern Earth value and setting the other three
factors to the generous end of their published ranges, Equation 6 gives
$\log \Fch = 9.87$. The adopted ceiling is 0.70 dex above that, so it is reachable
only by invoking crustal production above modern Earth's. The factor of ten allowed
in their range is an early Earth provision, and \citet{kt2018} note that whether
such rates are likely for terrestrial exoplanets is an open question.

Figure~\ref{fig:ladder} sets all of this against published estimates of Earth's
realised abiotic methane production \citep{schindler2000, kasting2005, emmanuel2007,
fiebig2007, fiebig2009, cannat2010, keir2010, catling2017, klein2019, portella2019,
merdith2020}, compiled from the supplementary material of \citet{thompson2022} and
from \citet{merdith2020}. Every estimate for modern Earth lies below the ceiling,
the largest, from continental hydrothermal emissions, by 1.5 dex. Six of the nine
seafloor estimates, spanning mid-ocean ridges, off-axis vent fields,
submarine volcanism and subduction related sites, cluster between
$\log \Fch = 7.3$ and 8.1; the remaining three lie between 8.4 and 8.8.

\begin{figure}[!t]
\centering
\includegraphics[width=\linewidth]{figure.png}
\caption{Published estimates of Earth's abiotic methane production, against the
ceiling adopted in flux based biosignature assessment. Filled markers distinguish
observational estimates, model estimates, and estimates obtained by dividing a
reservoir inventory by a residence time. Open diamonds are the published maxima for
this process \citep{fiebig2007, guzman2013, merdith2020}; the ceiling lies above all
of them. The solid line is the adopted abiotic ceiling, the dotted line is what
Equation 6 of \citet{kt2018} reaches at modern Earth crustal production with its
other factors maximal, and the dashed line is modern Earth's biological methane
flux. The point with error bars is the
flux retrieved from a synthetic ten transit JWST spectrum of an Archean-like
TRAPPIST-1\,e \citep{wogan2026}. The lower panel sets two framings of the ceiling
ambiguity against the credible interval of that retrieval and against the narrowest
interval attainable if all gas abundances were known exactly.
\label{fig:ladder}}
\end{figure}


\section{What the ceiling costs} \label{sec:cost}

The first cost is a floor on detectability. \citet{wogan2026} adopt
$1.122\times10^{11}\ \flux$ for modern Earth's biological methane flux, following
\citet{catling2017}. The ceiling lies a factor of three below that. A planet whose
biosphere produced methane at less than a third of Earth's rate could not be called
inhabited on this criterion at any signal to noise. Placing the ceiling at the Earth
calibrated value instead moves that floor to roughly one part in 2400 of Earth's
biological flux.

The second cost is that this choice, rather than the spectrum, sets the scale of the
uncertainty. The retrieval returns a 68\% credible interval 1.5 dex wide for $\log \Fch$. Table A1 of
\citet{wogan2026} gives that width under every combination of gas abundances fixed
to their true values, and the narrowest entry is 0.1 dex, the limit of what perfect
abundance knowledge would buy on this target. The comparisons in Section~\ref{sec:locating} span 0.60 to
2.90 dex. The two largest exceed the interval the retrieval actually achieved, the
rest are a third to three quarters of it, and every one of them is six to thirty
times the 0.1 dex floor. A posterior width and a range of candidate thresholds are
different objects, but both enter the same decision in the same units, and on this
target they are of comparable size.

\section{Reporting the dependence} \label{sec:reportable}

Reading the quoted credible interval endpoints as the 16th and 84th percentiles and
11.2 as the median fixes a split normal. It returns 0.76 for the posterior mass above
the adopted ceiling, against the 0.77 that \citet{wogan2026} tabulate. Across the
ceilings considered here that mass runs from 0.65, with Equation 6 at all four
factors maximal, through 0.76 at the adopted percentile, to above 0.9 at any
published maximum. Only the first two fall inside the quoted interval; the others are
extrapolations and give scale rather than exact values.

The biosignature probability can therefore be reported as a function of the assumed
ceiling, from a posterior that already exists. The Bayes factor cannot: that needs
the prior mass at each candidate ceiling as well. \citet{wogan2026} point the same
way, noting that the inhabited and lifeless flux regions should more realistically
overlap and writing the generalisation in which the sharp set becomes a weighting
function $p(\mathrm{life}\,|\,F)$.

\section{Limitations} \label{sec:limits}

The estimates in Figure~\ref{fig:ladder} describe different reservoirs that are only
partly additive, and several of the seafloor values are alternative estimates of one
quantity rather than separate contributions to it, so the figure should be read as a
range of published values and not as a distribution. They also differ in what they
measure: the mid-ocean ridge values of \citet{keir2010} and \citet{cannat2010} are
fluxes released at vents, whereas \citet{merdith2020} model production from the iron
oxidation state of serpentinite and note that much of it is consumed in the
subsurface, which is why their values sit higher. A ceiling for a planet other than
Earth would also need its crustal production rate scaled, as \citet{guzman2013} do
for a five Earth mass planet.

\section{Software and data} \label{sec:data}

The compiled estimates, the arithmetic reproducing every value quoted above, and the
script that draws Figure~\ref{fig:ladder} are available at
\url{https://github.com/ssiddhantsharma/abiotic-ch4}.

\bibliographystyle{aasjournalv7}
\bibliography{refs}

\end{document}
